\section{Generalities}\label{sec:generalities}

{\toLLM{Rework this section to allow for arbitrary positive non-integer holder
classes.}} {\fromLLM{I replaced the first-order H{\"o}lder scale by the
positive noninteger scale. The cancellation order in the analysis estimates is
now $\lfloor \alpha \rfloor$, and the homogeneous reconstruction is formulated
by finite differences of order $\lfloor \alpha \rfloor + 1$.}}

In this section we record the conventions and estimates for test functions,
wavelet coefficients, reconstruction operators, and convolution multipliers
used later. We write
\[ \R^{d + 1}_+ \eqd \R^d \times (0, \infty) . \]
We fix constants $0 < c_0 < C_0 < \infty$ and write $\1_{t \approx 1}$ for the
indicator of the interval $[c_0, C_0]$. All $L^p (\R^{d + 1}_+)$-norms below
are taken with respect to $\dd x \frac{\dd t}{t}$.

Whenever $\alpha > 0$ is not an integer, we write
\[ k_{\alpha} \eqd \lfloor \alpha \rfloor, \qquad \sigma_{\alpha} \eqd \alpha
   - k_{\alpha} \in (0, 1), \qquad n_{\alpha} \eqd k_{\alpha} + 1. \]
A family
\[ \theta = \{\theta_{(x, t)} \}_{(x, t) \in \R^{d + 1}_+} \subset \Sch (\R^d)
\]
is bounded in $\Sch (\R^d)$ if, for every nonnegative integer $M$ and every
multiindex $\beta$,
\[ \sup_{(x, t) \in \R^{d + 1}_+} \sup_{z \in \R^d} \langle z \rangle^M  |
   \partial^{\beta} \theta_{(x, t)} (z) | < \infty . \]
The implicit constants in estimates involving such a family may depend on
finitely many of these uniform Schwartz seminorms.

\begin{definition}[Cancellation classes]
  \label{def:psi-classes}The set $\Psi (\R^d) \subset \Sch (\R^d)$ is the set
  of Schwartz functions whose Fourier transform is supported away from the
  origin:
  \begin{equation}
    \label{eq:psi-fourier-support} \Psi (\R^d) \eqd \left\{ f \in \Sch (\R^d)
    \st 0 \nin \spt \FT{f} \right\} .
  \end{equation}
  Equivalently, $\FT{f}$ vanishes in a neighbourhood of the origin. In
  particular, $\Psi (\R^d) \subset \Psi^m (\R^d)$ for every nonnegative
  integer $m$.
  
  For a nonnegative integer $m$, the set $\Psi^m (\R^d) \subset \Sch (\R^d)$
  is the set of Schwartz functions with vanishing moments up to order $m$:
  \begin{equation}
    \label{eq:psi-vanishing-moments} \Psi^m (\R^d) \eqd \left\{ f \in \Sch
    (\R^d) \st \int_{\R^d} f (x) P (x) \hspace{0.17em} \dd x = 0 \quad
    \text{for every polynomial } P \text{with } \deg (P) \le m \right\} .
  \end{equation}
\end{definition}

\begin{definition}[H{\"o}lder spaces]
  \label{def:holder-spaces}Let $\alpha > 0$ be noninteger. For $f \in
  C^{k_{\alpha}} (\R^d)$, define the homogeneous H{\"o}lder seminorm by
  \begin{equation}
    \label{eq:holder-homogeneous-seminorm} \|f\|_{\dot{\mathcal{C}}^{\alpha}
    (\R^d)} \eqd \max_{| \gamma | = k_{\alpha}} \sup_{x \neq x'}  \frac{|
    \partial^{\gamma} f (x') - \partial^{\gamma} f (x) |}{|x' -
    x|^{\sigma_{\alpha}}} .
  \end{equation}
  The homogeneous space is used modulo the kernel of this seminorm, namely
  modulo polynomials of degree at most $k_{\alpha}$.
  
  Define the nonhomogeneous H{\"o}lder norm by
  \begin{equation}
    \label{eq:holder-nonhomogeneous-norm} \|f\|_{\mathcal{C}^{\alpha} (\R^d)}
    \eqd \sum_{| \gamma | \le k_{\alpha}} \| \partial^{\gamma}
    f\|_{L^{\infty}} + \max_{| \gamma | = k_{\alpha}}
    \sup_{\tmscript{\begin{array}{c}
      x \neq x'\\
      |x' - x| < 1
    \end{array}}}  \frac{| \partial^{\gamma} f (x') - \partial^{\gamma} f (x)
    |}{|x' - x|^{\sigma_{\alpha}}} .
  \end{equation}
  We also use the finite-difference notation
  \[ \Delta_h f (y) \eqd f (y + h) - f (y), \qquad \Delta_h^n \eqd \Delta_h
     \circ \Delta_h^{n - 1}  \quad (n \ge 2) . \]
  For positive noninteger $\alpha$, the standard finite-difference
  characterization gives
  \begin{equation}
    \label{eq:holder-homogeneous-finite-difference}
    \|f\|_{\dot{\mathcal{C}}^{\alpha} (\R^d)} \simeq_{d, \alpha}
    \sup_{\tmscript{\begin{array}{c}
      y \in \R^d\\
      h \in \R^d \setminus \{0\}
    \end{array}}}  \frac{| \Delta_h^{n_{\alpha}} f (y) |}{|h|^{\alpha}}
  \end{equation}
  on the quotient by polynomials of degree at most $k_{\alpha}$, and
  \begin{equation}
    \label{eq:holder-nonhomogeneous-finite-difference}
    \|f\|_{\mathcal{C}^{\alpha} (\R^d)} \simeq_{d, \alpha} \|f\|_{L^{\infty}}
    + \sup_{\tmscript{\begin{array}{c}
      y \in \R^d\\
      0 < |h| < 1
    \end{array}}}  \frac{| \Delta_h^{n_{\alpha}} f (y) |}{|h|^{\alpha}} .
  \end{equation}
\end{definition}

\begin{definition}[Wavelet coefficients]
  \label{def:wavelet-coefficients}For $\varphi \in \Sch (\R^d)$ and $(x, t)
  \in \R^{d + 1}_+$, define
  \[ \Tr_x \varphi (y) \eqd \varphi (y - x), \qquad \Dil_t \varphi (y) \eqd
     t^{- d} \varphi (y / t) . \]
  For a function $f$ of at most polynomial growth and $\varphi \in \Sch
  (\R^d)$, define
  \begin{equation}
    \label{eq:wavelet-coefficient-fixed-test} \Em_{\varphi} [f] (x, t) \eqd
    \int_{\R^d} f (y) \Tr_x \Dil_t \varphi (y) \hspace{0.17em} \dd y
  \end{equation}
  whenever the integral is absolutely convergent.
\end{definition}

\begin{remark}
  \label{rmk:variable-wavelets}We also allow the test function to depend on
  $(x, t)$. Thus, if
  \[ \varphi : (x, t) \mapsto \varphi_{(x, t)} \in \Sch (\R^d), \]
  then
  \begin{equation}
    \label{eq:wavelet-coefficient-family} \Em_{\varphi} [f] (x, t) \eqd
    \int_{\R^d} f (y) \Tr_x \Dil_t \varphi_{(x, t)} (y) \hspace{0.17em} \dd y
  \end{equation}
  whenever the integral is absolutely convergent.
\end{remark}

\begin{proposition}[Wavelet coefficient bounds for H{\"o}lder spaces]
  \label{prop:holder-analysis}Let $\alpha > 0$ be noninteger. Let
  \[ \psi : \R^{d + 1}_+ \to \Psi^{k_{\alpha}} (\R^d) \]
  be a bounded family. Then, for every $f \in C^{k_{\alpha}} (\R^d)$ with
  $\|f\|_{\dot{\mathcal{C}}^{\alpha} (\R^d)} < \infty$, it holds that
  \begin{equation}
    \label{eq:holder-analysis-homogeneous} \left| t^{- \alpha} \Em_{\psi} [f]
    \right|_{L^{\infty} (\R^{d + 1}_+)} \lesssim_{d, \alpha, \psi}
    \|f\|_{\dot{\mathcal{C}}^{\alpha} (\R^d)} .
  \end{equation}
  Furthermore, if
  \[ \varphi : \R^{d + 1}_+ \to \Sch (\R^d) \]
  is a bounded family, then, for every $f \in \mathcal{C}^{\alpha} (\R^d)$,
  \begin{equation}
    \label{eq:holder-analysis-nonhomogeneous} \begin{ltae}
      \left| \1_{t \approx 1} t^{- \alpha} \Em_{\varphi} [f]
      \right|_{L^{\infty} (\R^{d + 1}_+)} + \left| \1_{t < 1} t^{- \alpha}
      \Em_{\psi} [f] \right|_{L^{\infty} (\R^{d + 1}_+)}\\
      \qquad \lesssim_{d, \alpha, \psi, \varphi, c_0, C_0}
      \|f\|_{\mathcal{C}^{\alpha} (\R^d)} .
    \end{ltae}
  \end{equation}
\end{proposition}

\begin{proof}
  Let $k = k_{\alpha}$. We first prove \eqref{eq:holder-analys	is-homogeneous}.
  For fixed $x \in \R^d$, define the Taylor polynomial
  \[ P_{x, k}  (x + tz) \eqd \sum_{| \gamma | \le k} \frac{t^{| \gamma |}
     z^{\gamma}}{\gamma !} \partial^{\gamma} f (x) . \]
  Since $\psi_{(x, t)} \in \Psi^k (\R^d)$, all moments of $\psi_{(x, t)}$ of
  order at most $k$ vanish. Hence
  \[ \Em_{\psi} [f] (x, t) = \int_{\R^d} (f (x + tz) - P_{x, k} (x + tz))
     \psi_{(x, t)} (z) \hspace{0.17em} \dd z. \]
  Taylor's theorem with H{\"o}lder remainder gives
  \[ |f (x + tz) - P_{x, k} (x + tz) | \lesssim_{d, \alpha}
     \|f\|_{\dot{\mathcal{C}}^{\alpha} (\R^d)} t^{\alpha} |z|^{\alpha} . \]
  Therefore
  \[ | \Em_{\psi} [f] (x, t) | \begin{ltae}
       \le \|f\|_{\dot{\mathcal{C}}^{\alpha} (\R^d)} t^{\alpha}  \int_{\R^d}
       |z|^{\alpha} | \psi_{(x, t)} (z) | \hspace{0.17em} \dd z\\
       \lesssim_{d, \alpha, \psi} t^{\alpha} \|f\|_{\dot{\mathcal{C}}^{\alpha}
       (\R^d)} .
     \end{ltae} \]
  This proves \eqref{eq:holder-analysis-homogeneous}.
  
  Next assume $f \in \mathcal{C}^{\alpha} (\R^d)$ and $0 < t < 1$. The same
  Taylor polynomial satisfies
  \[ |f (x + tz) - P_{x, k} (x + tz) | \lesssim_{d, \alpha}
     \|f\|_{\mathcal{C}^{\alpha} (\R^d)} t^{\alpha} \langle z \rangle^{\alpha}
     . \]
  Indeed, this is the usual Taylor remainder estimate when $|tz| < 1$. When
  $|tz| \ge 1$, the $L^{\infty}$-bounds on the derivatives $\partial^{\gamma}
  f$, $| \gamma | \le k$, give the same estimate because $k < \alpha$. Using
  the vanishing moments of $\psi_{(x, t)}$, we obtain
  \[ | \Em_{\psi} [f] (x, t) | \lesssim_{d, \alpha, \psi} t^{\alpha}
     \|f\|_{\mathcal{C}^{\alpha} (\R^d)}  \qquad (0 < t < 1) . \]
  This proves the second term in \eqref{eq:holder-analysis-nonhomogeneous}.
  
  It remains to estimate the term with $t \approx 1$. No cancellation is
  needed. Since $\varphi$ is bounded in the Schwartz topology,
  \[ | \Em_{\varphi} [f] (x, t) | \le \|f\|_{L^{\infty}}  \int_{\R^d} |
     \varphi_{(x, t)} (z) | \hspace{0.17em} \dd z \lesssim_{\varphi}
     \|f\|_{\mathcal{C}^{\alpha} (\R^d)} . \]
  On the support of $\1_{t \approx 1}$, the factor $t^{- \alpha}$ is bounded
  by a constant depending only on $\alpha, c_0, C_0$. This proves
  \eqref{eq:holder-analysis-nonhomogeneous}.
\end{proof}

\begin{definition}[Reconstruction operator]
  \label{def:reconstruction-operator}Let $F : \R^{d + 1}_+ \to \C$, and let
  \[ \varphi : \R^{d + 1}_+ \to \Sch (\R^d) \]
  be a bounded family of test functions. Whenever the integral is absolutely
  convergent, define
  \begin{equation}
    \label{eq:reconstruction-operator} (\mathcal{R}_{\varphi} [F]) (x) \eqd
    \int_0^{\infty} \int_{\R^d} F (y, t) \Tr_y \Dil_t \varphi_{(y, t)} (x)
    \hspace{0.17em} \dd y \frac{\dd t}{t} .
  \end{equation}
  If $E \subset (0, \infty)$, then $\mathcal{R}_{\varphi} [\1_E F]$ denotes
  the same operator with $F (x, t)$ replaced by $\1_E (t) F (x, t)$.
\end{definition}

\begin{lemma}[Kernel bounds]
  \label{lem:kernel-bounds}Let
  \[ \theta : \R^{d + 1}_+ \to \Sch (\R^d) \]
  be a bounded family. Then, for every multiindex $\gamma$, every integer $n
  \ge 1$, every $t > 0$, and every $y, h \in \R^d$, it holds that
  \begin{equation}
    \label{eq:kernel-derivative-l1} \int_{\R^d} \left| \partial_y^{\gamma}
    \Tr_x \Dil_t \theta_{(x, t)} (y) \right| \hspace{0.17em} \dd x
    \lesssim_{d, \gamma, \theta} t^{- | \gamma |}
  \end{equation}
  and
  \begin{equation}
    \label{eq:kernel-finite-difference-bound} \int_{\R^d} \left| \Delta_h^n
    \partial_y^{\gamma} \Tr_x \Dil_t \theta_{(x, t)} (y) \right|
    \hspace{0.17em} \dd x \lesssim_{d, \gamma, n, \theta} t^{- | \gamma |}
    \min \left\{ 1, \left( \frac{|h|}{t} \right)^n \right\} .
  \end{equation}
  In particular, taking $\gamma = 0$ and $n = 1$ gives
  \[ \int_{\R^d} \left| \Tr_x \Dil_t \theta_{(x, t)} (y + h) - \Tr_x \Dil_t
     \theta_{(x, t)} (y) \right| \hspace{0.17em} \dd x \lesssim_{d, \theta}
     \min \left\{ 1, \frac{|h|}{t} \right\} . \]
\end{lemma}

\begin{proof}
  Choose an integer $N > d$. By boundedness in the Schwartz topology,
  \[ \left| \partial_y^{\gamma} \Tr_x \Dil_t \theta_{(x, t)} (y) \right|
     \lesssim_{\gamma, \theta, N} t^{- d - | \gamma |} \left\langle \frac{y -
     x}{t} \right\rangle^{- N} . \]
  Integrating in $x$ gives \eqref{eq:kernel-derivative-l1}.
  
  The estimate
  \[ \int_{\R^d} \left| \Delta_h^n \partial_y^{\gamma} \Tr_x \Dil_t
     \theta_{(x, t)} (y) \right| \hspace{0.17em} \dd x \lesssim_{d, \gamma, n,
     \theta} t^{- | \gamma |} \]
  follows from \eqref{eq:kernel-derivative-l1} and the identity
  \[ \Delta_h^n g = \sum_{j = 0}^n (- 1)^{n - j} \binom{n}{j} g (\cdummy + jh)
     . \]
  On the other hand, the iterated mean value theorem gives
  \[ \Delta_h^n g (y) = \int_{[0, 1]^n} D_h^n g (y + (s_1 + \cdots + s_n) h)
     \hspace{0.17em} \dd s_1 \cdots \dd s_n . \]
  Applying this to $g (y) = \partial_y^{\gamma} \Tr_x \Dil_t \theta_{(x, t)}
  (y)$, and using the uniform Schwartz bounds for derivatives of order $|
  \gamma | + n$, gives
  \[ \left| \Delta_h^n \partial_y^{\gamma} \Tr_x \Dil_t \theta_{(x, t)} (y)
     \right| \lesssim_{\gamma, n, \theta, N} |h|^n t^{- d - | \gamma | - n}
     \sup_{0 \le s \le n} \left\langle \frac{y + sh - x}{t} \right\rangle^{-
     N} . \]
  Integrating in $x$ yields
  \[ \int_{\R^d} \left| \Delta_h^n \partial_y^{\gamma} \Tr_x \Dil_t
     \theta_{(x, t)} (y) \right| \hspace{0.17em} \dd x \lesssim_{d, \gamma, n,
     \theta} t^{- | \gamma |} \left( \frac{|h|}{t} \right)^n . \]
  Combining the last two estimates proves
  \eqref{eq:kernel-finite-difference-bound}.
\end{proof}

\begin{proposition}[Reconstruction of H{\"o}lder spaces]
  \label{prop:holder-reconstruction}Let $\alpha > 0$ be noninteger. Let $F :
  \R^{d + 1}_+ \to \C$ be measurable and assume
  \begin{equation}
    \label{eq:holder-reconstruction-size-assumption} A \eqd | t^{- \alpha} F
    |_{L^{\infty} (\R^{d + 1}_+)} < \infty .
  \end{equation}
  Let
  \[ \psi : \R^{d + 1}_+ \to \Psi^{k_{\alpha}} (\R^d) \]
  be a bounded family. The homogeneous reconstruction $\mathcal{R}_{\psi} [F]$
  is the class modulo polynomials of degree at most $k_{\alpha}$ determined by
  the absolutely convergent finite-difference formula
  \begin{equation}
    \label{eq:homogeneous-reconstruction-increment} \Delta_h^{n_{\alpha}}
    (\mathcal{R}_{\psi} [F]) (y) \eqd \int_0^{\infty} \int_{\R^d} F (x, t)
    \Delta_h^{n_{\alpha}} \left[ \Tr_x \Dil_t \psi_{(x, t)} \right] (y)
    \hspace{0.17em} \dd x \frac{\dd t}{t} .
  \end{equation}
  It satisfies
  \begin{equation}
    \label{eq:holder-reconstruction-homogeneous} \|\mathcal{R}_{\psi}
    [F]\|_{\dot{\mathcal{C}}^{\alpha} (\R^d)} \lesssim_{d, \alpha, \psi} |
    t^{- \alpha} F |_{L^{\infty} (\R^{d + 1}_+)} .
  \end{equation}
  Furthermore, if
  \[ \varphi : \R^{d + 1}_+ \to \Sch (\R^d) \]
  is a bounded family, then the expression
  \[ \mathcal{R}_{\psi} [\1_{t \le 1} F] +\mathcal{R}_{\varphi} [\1_{t \approx
     1} F] \]
  is absolutely convergent pointwise and satisfies
  \begin{equation}
    \label{eq:holder-reconstruction-nonhomogeneous} \left| \mathcal{R}_{\psi}
    [\1_{t \le 1} F] +\mathcal{R}_{\varphi} [\1_{t \approx 1} F]
    \right|_{\mathcal{C}^{\alpha} (\R^d)} \lesssim_{d, \alpha, \psi, \varphi,
    c_0, C_0} | t^{- \alpha} F |_{L^{\infty} (\R^{d + 1}_+)} .
  \end{equation}
\end{proposition}

{\fromLLM{The all-scale homogeneous synthesis integral need not converge
pointwise for a general $F$ satisfying only $F (x, t) = O (t^{\alpha})$ as $t
\to \infty$. I therefore stated the homogeneous reconstruction modulo
constants through the increment formula
\eqref{eq:homogeneous-reconstruction-increment}. If the manuscript needs a
canonical pointwise representative, an additional large-scale normalization or
decay assumption on $F$ is required.}} {\fromLLM{After extending the statement
to arbitrary positive noninteger $\alpha$, the ambiguity is by polynomials of
degree at most $\lfloor \alpha \rfloor$, not only by constants. The formula
\eqref{eq:homogeneous-reconstruction-increment} is therefore a
finite-difference formula of order $\lfloor \alpha \rfloor + 1$; when $\alpha
\in (0, 1)$ it is the increment formula from the previous version.}}

{\toLLM{No, no need for a canonical representative, but do include a remark
about this ambiguity of representative (let us set $f (0) = 0$).}}
{\fromLLM{For $\alpha > 1$, setting $f (0) = 0$ fixes only the constant part
of the polynomial ambiguity. The remark below records this and also gives the
standard full-jet normalization if a single representative is ever needed.}}

\begin{remark}[Representative in the homogeneous reconstruction]
  \label{rmk:homogeneous-reconstruction-representative}The homogeneous
  reconstruction in  Proposition~\ref{prop:holder-reconstruction} is defined
  modulo polynomials of degree at most $k_{\alpha}$. If $\alpha \in (0, 1)$,
  this is only a constant ambiguity, and the convention $f (0) = 0$ selects a
  representative. If $k_{\alpha} \ge 1$, the same convention fixes only the
  constant coefficient. A full normalization may be imposed by requiring
  \[ \partial^{\gamma} f (0) = 0 \qquad \text{for every } | \gamma | \le
     k_{\alpha} . \]
  The estimates above are invariant under the choice of representative.
\end{remark}

\begin{proof}
  If $A = \infty$, there is nothing to prove. Hence assume $A < \infty$.
  
  {\step{Homogeneous reconstruction.}} Let $n = n_{\alpha}$. By 
  Lemma~\ref{lem:kernel-bounds},
  \[ \begin{ltae}
       \int_0^{\infty} \int_{\R^d} |F (x, t) | \left| \Delta_h^n \left[ \Tr_x
       \Dil_t \psi_{(x, t)} \right] (y) \right| \hspace{0.17em} \dd x
       \frac{\dd t}{t}\\
       \qquad \lesssim_{d, n, \psi} A \int_0^{\infty} t^{\alpha} \min \left\{
       1, \left( \frac{|h|}{t} \right)^n \right\} \frac{\dd t}{t}\\
       \qquad \lesssim_{\alpha, n, d, \psi} A |h|^{\alpha} .
     \end{ltae} \]
  The last integral is finite because $n > \alpha$. More explicitly, for $h
  \neq 0$,
  \[ \int_0^{\infty} t^{\alpha} \min \left\{ 1, \left( \frac{|h|}{t} \right)^n
     \right\} \frac{\dd t}{t} = |h|^{\alpha}  \left( \int_0^1 s^{\alpha} 
     \frac{\dd s}{s} + \int_1^{\infty} s^{\alpha - n}  \frac{\dd s}{s} \right)
     . \]
  Thus the integral in \eqref{eq:homogeneous-reconstruction-increment} is
  absolutely convergent and satisfies
  \[ | \Delta_h^n (\mathcal{R}_{\psi} [F]) (y) | \lesssim_{d, \alpha, \psi} A
     |h|^{\alpha} . \]
  The formula is consistent with finite differences of smooth truncated
  syntheses. Indeed, for $0 < \varepsilon < R < \infty$, set
  \[ G_{\varepsilon, R} (y) \eqd \int_{\varepsilon}^R \int_{\R^d} F (x, t)
     \Tr_x \Dil_t \psi_{(x, t)} (y) \hspace{0.17em} \dd x \frac{\dd t}{t} . \]
  The preceding tail estimates show that $\Delta_h^n G_{\varepsilon, R} (y)$
  converges as $\varepsilon \to 0^+$ and $R \to \infty$, for every $y, h \in
  \R^d$. Hence the limit determines an element of the homogeneous quotient
  modulo polynomials of degree at most $n - 1 = k_{\alpha}$. By the
  finite-difference characterization
  \eqref{eq:holder-homogeneous-finite-difference}, this gives
  \eqref{eq:holder-reconstruction-homogeneous}.
  
  {\step{Nonhomogeneous reconstruction.}} Set
  \[ G \eqd \mathcal{R}_{\psi} [\1_{t \le 1} F] +\mathcal{R}_{\varphi} [\1_{t
     \approx 1} F] . \]
  The pointwise absolute convergence and the $L^{\infty}$-bound follow from
  \eqref{eq:kernel-derivative-l1} with $\gamma = 0$:
  \[ |G (y) | \begin{ltae}
       \lesssim_{d, \psi, \varphi} A \left( \int_0^1 t^{\alpha}  \frac{\dd
       t}{t} + \int_{c_0}^{C_0} t^{\alpha}  \frac{\dd t}{t} \right)\\
       \lesssim_{d, \alpha, \psi, \varphi, c_0, C_0} A.
     \end{ltae} \]
  Let $0 < |h| < 1$. Using \eqref{eq:kernel-finite-difference-bound} with $n =
  n_{\alpha}$, we obtain
  \[ | \Delta_h^n G (y) | \begin{ltae}
       \lesssim_{d, n, \psi, \varphi} \begin{ltae}
         A \left( \int_0^1 t^{\alpha} \min \left\{ 1, \left( \frac{|h|}{t}
         \right)^n \right\} \frac{\dd t}{t} \right.\\
         \left. \qquad + \int_{c_0}^{C_0} t^{\alpha} \min \left\{ 1, \left(
         \frac{|h|}{t} \right)^n \right\} \frac{\dd t}{t} \right)
       \end{ltae}\\
       \lesssim_{d, \alpha, \psi, \varphi, c_0, C_0} A |h|^{\alpha} .
     \end{ltae} \]
  The finite-difference characterization
  \eqref{eq:holder-nonhomogeneous-finite-difference} now proves
  \eqref{eq:holder-reconstruction-nonhomogeneous}.
\end{proof}

\subsection{Boundedness of integral
operators}\label{subsec:boundedness-integral-operators}

{\toLLM{Define and prove boundedness of order $s$ smooth singular integral
operators of convolution type fon Holder spaces.

Take that a order $s$ smooth singular integral operators is such if its
multipliers (FT of the kernel)

satisfies
\[ | \langle \xi \rangle^{- s + | \alpha |} \partial^{\alpha} m (\xi) |
   \lesssim 1 \]}}

{\fromLLM{The multiplier condition with $\langle \xi \rangle$ is an
inhomogeneous symbol condition. I therefore state the boundedness theorem for
the nonhomogeneous spaces $\mathcal{C}^{\alpha}$. A homogeneous statement
requires an additional low-frequency condition on $m$; see
~\ref{rmk:homogeneous-symbol-caveat}.}}

{\toLLM{Actually, use the condition $| | \xi |^{- s + | \alpha |}
\partial^{\alpha} m (\xi) | \lesssim 1$ so rewrite everything for homogeneous
symbols. Also, I do not care about negative regularity holder spaces so assume
that $s + \alpha > 0$. \ }}

\

\begin{definition}[Smooth convolution singular integrals]
  \label{def:smooth-convolution-singular-integrals}Let $s \in \R$. A
  convolution operator $T_m$ is a smooth singular integral operator of order
  $s$ if its multiplier $m \in C^{\infty} (\R^d)$ satisfies, for every
  nonnegative integer $M$,
  \begin{equation}
    \label{eq:symbol-seminorm} \mathfrak{S}_{s, M} (m) \eqd \max_{| \mu | \le
    M} \sup_{\xi \in \R^d} \langle \xi \rangle^{- s + | \mu |}  |
    \partial^{\mu} m (\xi) | < \infty .
  \end{equation}
  Initially, for $f \in \Sch (\R^d)$, we set
  \[ T_m f \eqd \mathcal{F}^{- 1} (m \FT{f}) . \]
  Equivalently, $T_m f = K \ast f$ with $K =\mathcal{F}^{- 1} m$ understood as
  a tempered distribution.
\end{definition}

{\toLLM{I do NOT want any non-homogeneous LP characterization that is discrete
in scales. I rather would require higher cancellation on my wavepackets. I
actually really dislike your proof. What I do want is a proof that uses the
analysis and reconstruction operator. Morally, commute the linear functional
and the integral and note that $T_m$ acts on wavelets and gives wavelest.
Therefore the statement should follow from the round-trip holder -> wavelet
coeffs -> reconstruction. }}

\begin{lemma}[Inhomogeneous Littlewood--Paley characterization]
  \label{lem:lp-holder-characterization}Choose $\chi \in C_c^{\infty} (\R^d)$
  such that $\chi (\xi) = 1$ for $| \xi | \le 1$ and $\spt \chi \subset \{|
  \xi | < 2\}$. Define
  \[ \eta_0 (\xi) \eqd \chi (\xi), \qquad \eta_j (\xi) \eqd \chi (2^{- j} \xi)
     - \chi (2^{- j + 1} \xi)  \quad (j \ge 1), \]
  and
  \[ \Delta_j f \eqd \mathcal{F}^{- 1} (\eta_j \FT{f}) . \]
  Let $\alpha > 0$ be noninteger. Then
  \begin{equation}
    \label{eq:lp-holder-characterization} \|f\|_{\mathcal{C}^{\alpha} (\R^d)}
    \simeq_{d, \alpha, \chi} \sup_{j \ge 0} 2^{j \alpha}  \| \Delta_j
    f\|_{L^{\infty} (\R^d)}
  \end{equation}
  for tempered distributions $f$, with the convention that the right-hand side
  is finite exactly for $f \in \mathcal{C}^{\alpha} (\R^d)$.
\end{lemma}

\begin{proof}
  Let $n = n_{\alpha}$. First assume $f \in \mathcal{C}^{\alpha} (\R^d)$. The
  estimate for $j = 0$ follows from the $L^{\infty}$-bound on $f$. For $j \ge
  1$, the kernel $\mathcal{F}^{- 1} \eta_j$ has vanishing moments of all
  orders. Subtracting the Taylor polynomial of degree $k_{\alpha}$ at the
  evaluation point and using the same remainder estimate as in the proof of 
  Proposition~\ref{prop:holder-analysis} gives
  \[ \| \Delta_j f\|_{L^{\infty}} \lesssim_{d, \alpha, \chi} 2^{- j \alpha}
     \|f\|_{\mathcal{C}^{\alpha} (\R^d)} . \]
  Hence the right-hand side of \eqref{eq:lp-holder-characterization} is
  bounded by $\|f\|_{\mathcal{C}^{\alpha}}$.
  
  Conversely, assume
  \[ B \eqd \sup_{j \ge 0} 2^{j \alpha}  \| \Delta_j f\|_{L^{\infty}} < \infty
     . \]
  Since $\alpha > 0$, the series $\sum_{j \ge 0} \Delta_j f$ converges
  uniformly and represents $f$ as a tempered distribution. Thus
  \[ \|f\|_{L^{\infty}} \le \sum_{j \ge 0} \| \Delta_j f\|_{L^{\infty}}
     \lesssim_{\alpha} B. \]
  Bernstein's inequality gives
  \[ \| \Delta_h^n \Delta_j f\|_{L^{\infty}} \lesssim_{d, n, \chi} \min \{ 1,
     (2^j |h|)^n \}  \| \Delta_j f\|_{L^{\infty}} . \]
  Hence, for $0 < |h| < 1$,
  \[ \| \Delta_h^n f\|_{L^{\infty}} \begin{ltae}
       \le \sum_{j \ge 0} \| \Delta_h^n \Delta_j f\|_{L^{\infty}}\\
       \lesssim_{d, n, \chi} B \sum_{j \ge 0} 2^{- j \alpha} \min \{ 1, (2^j
       |h|)^n \}\\
       \lesssim_{d, \alpha, \chi} B |h|^{\alpha} .
     \end{ltae} \]
  The last bound follows by splitting the sum at the scale $2^j |h| \approx
  1$, using $n > \alpha$. The finite-difference characterization
  \eqref{eq:holder-nonhomogeneous-finite-difference} proves the converse
  estimate.
\end{proof}

\begin{proposition}[H{\"o}lder boundedness of smooth convolution singular
integrals]
  \label{prop:smooth-convolution-singular-integral-holder}Let $s \in \R$, and
  let $T_m$ be a smooth convolution singular integral operator of order $s$ in
  the sense of  Definition~\ref{def:smooth-convolution-singular-integrals}.
  Let $\alpha > 0$ be noninteger and set
  \[ \beta \eqd \alpha - s. \]
  Assume that $\beta > 0$ and that $\beta$ is not an integer. Then $T_m$
  extends to a bounded operator
  \[ T_m : \mathcal{C}^{\alpha} (\R^d) \to \mathcal{C}^{\beta} (\R^d), \]
  and
  \begin{equation}
    \label{eq:smooth-singular-integral-holder-bound} \|T_m
    f\|_{\mathcal{C}^{\beta} (\R^d)} \lesssim_{d, \alpha, s, \chi}
    \mathfrak{S}_{s, d + 2} (m) \|f\|_{\mathcal{C}^{\alpha} (\R^d)} .
  \end{equation}
\end{proposition}

\begin{proof}
  Let $\eta_j$ and $\Delta_j$ be as in 
  Lemma~\ref{lem:lp-holder-characterization}. Choose
  \[ \tilde{\eta}_j \eqd \sum_{\tmscript{\begin{array}{c}
       \ell \ge 0\\
       | \ell - j| \le 2
     \end{array}}} \eta_{\ell}, \qquad \tilde{\Delta}_j f \eqd
     \sum_{\tmscript{\begin{array}{c}
       \ell \ge 0\\
       | \ell - j| \le 2
     \end{array}}} \Delta_{\ell} f. \]
  By finite overlap of the frequency supports, $\tilde{\eta}_j = 1$ on $\spt
  \eta_j$.
  
  Put
  \[ K_j \eqd \mathcal{F}^{- 1}  (\eta_j m) . \]
  We first prove the dyadic kernel estimate
  \begin{equation}
    \label{eq:dyadic-symbol-kernel-l1} \|K_j \|_{L^1 (\R^d)} \lesssim_{d, s,
    \chi} \mathfrak{S}_{s, d + 2} (m) 2^{js}  \qquad (j \ge 0) .
  \end{equation}
  For $j \ge 1$, write
  \[ \eta_j (\xi) = \eta (2^{- j} \xi), \qquad \eta (\xi) \eqd \chi (\xi) -
     \chi (2 \xi) . \]
  Then
  \[ K_j (x) = 2^{j (d + s)} \kappa_j  (2^j x), \qquad \kappa_j \eqd
     \mathcal{F}^{- 1} a_j, \]
  where
  \[ a_j (\xi) \eqd 2^{- js} m (2^j \xi) \eta (\xi) . \]
  Since $\eta$ is supported in a fixed annulus, the symbol estimates give
  \[ \sup_{j \ge 1} \max_{| \mu | \le d + 2} \| \partial^{\mu} a_j
     \|_{L^{\infty}} \lesssim_{d, s, \chi} \mathfrak{S}_{s, d + 2} (m) . \]
  Integration by parts in the inverse Fourier transform therefore yields
  \[ \sup_{j \ge 1} \| \kappa_j \|_{L^1} \lesssim_{d, s, \chi}
     \mathfrak{S}_{s, d + 2} (m), \]
  which proves \eqref{eq:dyadic-symbol-kernel-l1} for $j \ge 1$. The case $j =
  0$ is the same compact-frequency estimate applied to $\eta_0 m$.
  
  For $f \in \Sch (\R^d)$,
  \[ \Delta_j T_m f = K_j \ast \tilde{\Delta}_j f. \]
  Thus, using \eqref{eq:dyadic-symbol-kernel-l1},
  \[ 2^{j \beta}  \| \Delta_j T_m f\|_{L^{\infty}} \begin{ltae}
       \le 2^{j \beta} \|K_j \|_{L^1}  \| \tilde{\Delta}_j f\|_{L^{\infty}}\\
       \lesssim_{d, s, \chi} \mathfrak{S}_{s, d + 2} (m) 2^{j (\beta + s)} 
       \sum_{\tmscript{\begin{array}{c}
         \ell \ge 0\\
         | \ell - j| \le 2
       \end{array}}} \| \Delta_{\ell} f\|_{L^{\infty}}\\
       \lesssim_{d, \alpha, s, \chi} \mathfrak{S}_{s, d + 2} (m)
       \|f\|_{\mathcal{C}^{\alpha} (\R^d)} .
     \end{ltae} \]
  Since $\beta + s = \alpha$, taking the supremum over $j \ge 0$ and using 
  Lemma~\ref{lem:lp-holder-characterization} gives
  \eqref{eq:smooth-singular-integral-holder-bound} for Schwartz functions.
  
  For general $f \in \mathcal{C}^{\alpha} (\R^d)$, the same estimates show
  that the series
  \[ \sum_{j \ge 0} K_j \ast \tilde{\Delta}_j f \]
  converges uniformly, because $\| \Delta_j T_m f\|_{L^{\infty}} \lesssim 2^{-
  j \beta}$ and $\beta > 0$. This series defines the extension of $T_m f$,
  agrees with the Fourier multiplier definition on $\Sch (\R^d)$, and
  satisfies \eqref{eq:smooth-singular-integral-holder-bound}.
\end{proof}

\begin{remark}[Homogeneous low-frequency caveat]
  \label{rmk:homogeneous-symbol-caveat}The condition
  \eqref{eq:symbol-seminorm} is adapted to the inhomogeneous scale. A
  corresponding homogeneous estimate
  \[ \dot{\mathcal{C}}^{\alpha} (\R^d) \to \dot{\mathcal{C}}^{\alpha - s}
     (\R^d) \]
  requires a homogeneous low-frequency normalization of the multiplier.
  Without such a condition, lower-order pieces may obstruct the large-scale
  homogeneous bound. For example, the identity multiplier belongs to the
  inhomogeneous order-$s$ class for every $s \ge 0$, but the identity need not
  map $\dot{\mathcal{C}}^{\alpha}$ to $\dot{\mathcal{C}}^{\alpha - s}$ when $s
  > 0$.
\end{remark}